\section{Implementation Details}\label{app:implementation}

This appendix provides annotated code for the H\&M application and a diagnostics interpretation guide.

\subsection{Installation}

\begin{lstlisting}[language=bash, numbers=none]
pip install deep-inference
\end{lstlisting}

\noindent The package requires Python $\geq$ 3.8, PyTorch $\geq$ 1.12, NumPy, and scikit-learn (for ridge Lambda estimation). GPU support is automatic via PyTorch's CUDA backend.

\subsection{Data Preparation}

The data preparation script loads pre-trained embeddings, applies PCA to item embeddings, and constructs choice sets.

\begin{lstlisting}
import numpy as np
from sklearn.decomposition import PCA

# Load pre-trained two-tower embeddings
user_emb = np.load("artifacts/user_embeddings.npy")  # (5000, 64)
item_emb = np.load("artifacts/item_embeddings.npy")  # (90, 64)
prices = np.load("artifacts/ref_prices.npy")          # (90,)

# PCA on item embeddings: 64D -> 5D
pca = PCA(n_components=5)
item_pca = pca.fit_transform(item_emb)
print(f"Variance explained: {pca.explained_variance_ratio_.sum():.1%}")

# Build item attribute matrix: [log_price, pca_1, ..., pca_5]
log_prices = np.log(prices + 1e-6)
item_attrs = np.column_stack([log_prices, item_pca])  # (90, 6)

# For each purchase occasion, construct choice set
J = 20  # 1 chosen + 19 random alternatives
K = 6   # log_price + 5 PCA dims
rng = np.random.RandomState(42)

Y_list, T_list, X_list = [], [], []

for user_id, chosen_item_id in transactions:
    # Consumer embedding
    X_list.append(user_emb[user_id])

    # Pack chosen item first, then J-1 random alternatives
    t_row = np.zeros(J * K)
    t_row[:K] = item_attrs[chosen_item_id]

    others = np.setdiff1d(np.arange(len(item_attrs)), [chosen_item_id])
    sampled = rng.choice(others, J - 1, replace=False)
    for j, alt_id in enumerate(sampled):
        t_row[(j+1)*K : (j+2)*K] = item_attrs[alt_id]

    T_list.append(t_row)
    Y_list.append(0.0)  # Chosen is always index 0

Y = np.array(Y_list, dtype=np.float32)
T = np.array(T_list, dtype=np.float32)
X = np.array(X_list, dtype=np.float32)
np.save("data/Y.npy", Y)
np.save("data/T.npy", T)
np.save("data/X.npy", X)
\end{lstlisting}

\subsection{Inference Call}

The inference call is concise once the data is prepared:

\begin{lstlisting}
import torch
import numpy as np
from deep_inference import inference

# Load prepared data
Y = np.load("data/Y.npy")    # (n,)
T = np.load("data/T.npy")    # (n, 120)
X = np.load("data/X.npy")    # (n, 64)

# Custom MNL loss (5 lines)
J, K = 20, 6

def mnl_loss(y, t, theta):
    x = t.reshape(J, K)            # (20, 6) attribute matrix
    V = x @ theta                   # (20,) utilities
    return -V[0] + torch.logsumexp(V, dim=0)

# Target: average price sensitivity
def price_target(x, theta, t_tilde):
    return theta[0]

# Run inference
result = inference(
    Y=Y, T=T, X=X,
    loss=mnl_loss,
    theta_dim=K,                    # 6 parameters per consumer
    hessian_depends_on_theta=True,  # Softmax probs depend on theta
    hessian_depends_on_y=False,     # Fisher information Hessian
    target_fn=price_target,
    n_folds=50,                     # Cross-fitting folds
    epochs=300,                     # Training epochs
    patience=50,                    # CRITICAL for 3-way split
    hidden_dims=[64, 32],           # Network architecture
    lr=0.01,                        # Learning rate
    verbose=True,                   # Print progress
)

# Results
print(f"E[beta_price] = {result.mu_hat:.4f} +/- {result.se:.4f}")
print(f"95% CI: [{result.ci_lower:.4f}, {result.ci_upper:.4f}]")

# Diagnostics
print(f"Regime: {result.diagnostics.get('regime', 'C')}")
print(f"Lambda method: {result.diagnostics.get('lambda_method')}")
\end{lstlisting}

\subsection{Running All Parameters}

To estimate all $K$ parameters, loop over target indices:

\begin{lstlisting}
param_names = ["log_price", "pca_1", "pca_2", "pca_3", "pca_4", "pca_5"]

for k, name in enumerate(param_names):
    target_fn = lambda x, theta, t_tilde, idx=k: theta[idx]

    result = inference(
        Y=Y, T=T, X=X,
        loss=mnl_loss, theta_dim=K,
        hessian_depends_on_theta=True,
        target_fn=target_fn,
        n_folds=50, epochs=300, patience=50,
    )
    print(f"E[{name}] = {result.mu_hat:.4f} +/- {result.se:.4f}")
\end{lstlisting}

\subsection{Diagnostics Interpretation Guide}

\begin{enumerate}
    \item \textbf{Check the regime.} The output should say ``Regime C'' for our multinomial model. If it says Regime A or B, something is wrong with the Hessian detection.

    \item \textbf{Check Lambda eigenvalues.} If a warning about small eigenvalues appears, the model may be poorly identified. Consider: (a) reducing $d_\theta$ (more PCA compression), (b) increasing $n$, or (c) using a different Lambda method.

    \item \textbf{Check the correction ratio.} For multinomial logit, values of 70--90 are normal. For binary logit, values of 2--3 are normal. Very large values ($> 200$) suggest Lambda estimation problems.

    \item \textbf{Verify with simulation.} Before trusting results on real data, run the simulation study (\texttt{application/06\_simulation\_study.py}) with a DGP matching your application specification. If simulation coverage is near 95\%, you can be confident in the real-data results.

    \item \textbf{Compare IF to naive SE.} The IF SE should be larger than the naive SE. If it is smaller, something is wrong. The ratio depends on the model: 2--5$\times$ for binary logit, potentially larger for multinomial.
\end{enumerate}

\subsection{Companion Repository}

The complete code for the application is available at:

\begin{center}
    \texttt{https://github.com/rawatpranjal/deep-inference/tree/main/application}
\end{center}

\noindent The repository includes:
\begin{itemize}
    \item Pre-trained embedding artifacts (via symlink)
    \item Data preparation, inference, comparison, and counterfactual scripts
    \item Simulation study with coverage validation
    \item Figure generation for the paper
    \item The \texttt{deep-inference} package source code
\end{itemize}
